\documentclass{article}
\usepackage{pathfinder-readable}

\title{When Peer Comparisons Can Serve as Predictions}

\begin{document}
\maketitle

\begin{abstract}
This note reads Zhao and colleagues' survey of learning-augmented algorithms alongside a local paper on
rankings built from peer comparisons. We asked whether the latter's scores or rankings could be supplied to a
decision algorithm with a stated error and a stated cost of obtaining them. Peer agreements can pass the
paper's calibration check even when a shared error reverses a ranking; within a narrower shared-error model,
graph structure determines when order or numerical scores can still be recovered. A conditional sampling bound
prices the clean-model estimate, but no downstream decision guarantee follows. The verifier left the thread at
DRAFT because the argument supports these qualified findings and leaves the next interface theorem open.
\end{abstract}

\section{Introduction}

The survey by Zhao and colleagues, called Q here, studies algorithms that use fallible predictions while
retaining performance guarantees \cite{Q}. It distinguishes a numerical prediction from an ordering, asks how
prediction error affects a decision, and treats the cost of obtaining predictions as part of the system
question. It also explains why guarantees for separate parts of a system do not automatically give a guarantee
for the whole system.

The local paper \emph{From Truthful Peer Signals to Contribution Rankings}, called P here, studies how reports
from several observers can yield a numerical ranking \cite{P}. It assumes that observers compare the same
realised item, make independent symmetric errors conditional on that item, and have positive, stable
reliabilities. Agreements between observers then calibrate those reliabilities; the calibrated comparisons
determine Bradley--Terry scores when the item-comparison graph is connected. P also discusses using scores for
reward shaping. Its own letters Q and P name a different pair of papers.

The question pursued in this thread was whether P's output could be the prediction consumed by an algorithm of
the kind Q surveys. That requires a choice between an \emph{order} and numerical \emph{scores}, an error bound
for that choice, a count of comparison labels, and a rule for carrying the error into a particular decision. P
gives a population identification result under its assumptions. It does not itself give the finite-sample or
downstream rule.

\section{What was done}

The peers first checked which papers were actually supplied. They then translated P's population result into
Q's terms: a prediction must have a declared target, an error measure, and an acquisition cost. One peer
worked out how errors in measured agreements and edge means pass through P's reliability calibration,
Bradley--Terry inversion, and graph reconstruction. The other tested whether the agreements themselves show
that the intended target was recovered. The resulting four-vertex path example changed the focus from
estimation alone to identification under shared errors.

Next, the peers restricted the shared error to one positive attenuation rate common to all compared items.
They checked which comparison graphs preserve order without learning that rate, and which cycles can recover
the rate and hence the numerical scores. The triangle calculation was followed by a four-cycle calculation,
then by a six-cycle counterexample. Their derivations give the detailed calculations; the results below state
the assumptions needed to read them.

\section{Results}

Under P's clean model, let \(X_e\) be the binary result of comparison \(e\), and let \(Y_{er}\) be observer
\(r\)'s binary report on the same realised result. Write \(a_r>0\) for that observer's reliability.
Conditional independence and symmetric errors give
\[
q_{rs}:=\mathbb{E}[Y_{er}Y_{es}]=a_ra_s.
\]
Here \(q_{rs}\) is the agreement moment for observers \(r\) and \(s\). An observed triangle of observer pairs
fixes their positive reliabilities, and observed overlap paths extend the calibration. Dividing each edge's
report mean by a calibrated reliability recovers its latent comparison mean. P then uses a connected
comparison graph to reconstruct Bradley--Terry scores, up to adding the same constant to every score. This is
a conditional population statement, not an empirical test of its assumptions \cite{P}. The agreement step has
a precedent in single-coin crowdsourcing, as P notes \cite{P}.

The peers obtained a conditional sampling bound for this clean reconstruction. They assumed a positive lower
bound on reliability, edge probabilities bounded away from zero and one, bounded depth for the observer
calibration paths, and uniformly small errors in the measured moments. Their calculation bounded score error
by the moment error multiplied by factors that grow with calibration depth and with poor connectivity of the
comparison graph. With \(N\) independent repeats for each measured population, its sampling part decreases at
order \(N^{-1/2}\); a separate misspecification term may remain. If \(L_H\) is the set of measured
observer-overlap links and \(E\) the set of comparison edges, their direct collection plan uses at most
\(N(2|L_H|+|E|)\) observer labels. This is an upper bound for a specified estimator, with its derivation in
\texttt{emmy/working\_note.md}. It does not price those labels in a decision objective or establish a policy
guarantee.

The principal negative result concerns a shared error that all observers make on one item. Let \(B_e\) be a
shared binary sign, independent of the intended comparison \(X_e\), with the same positive mean \(b\) on every
edge. Let \(Z_{er}\) be an independent observer-specific sign of mean \(a_r\), and suppose each report is
\(Y_{er}=X_eB_eZ_{er}\). The shared sign cancels when two reports are multiplied. Thus their agreements still
obey \(q_{rs}=a_ra_s\), although the calibrated edge mean is
\[
w_{ij}=b\tanh\!\left(\frac{c_i-c_j}{2}\right),
\]
where \(c_i\) and \(c_j\) are the intended numerical scores and \(w_{ij}\) is the calibrated mean for their
comparison. Agreements alone therefore cannot certify that the inferred score targets the intended comparison.
This observation sharpens P's warning about correlated errors: the agreement factorisation can survive the
correlation. It does not contradict P's theorem, whose conditional independence premise excludes this shared
sign.

On a path of four items, the peers chose successive intended score differences \((4,-1.5,-1.5)\). The first
item's score exceeds the fourth's by \(1\). With \(b=1/2\), those same reports fit a clean Bradley--Terry
model whose first score falls below the fourth's by about \(0.265\). The two models give the same full
distribution of observed reports, even with unlimited repeated items. The construction and its
error-probability argument are in \texttt{ada/common\_mode\_counterexample.md}, and the second peer checked
the calculation in \texttt{emmy/working\_note.md}.

The graph matters differently for order and for numerical scores. Under the same constant positive
shared-error model, any two vertices joined by at most two comparison edges keep their relative order: the
edge transformation is odd and strictly increasing. Hence a connected comparison graph of diameter at most two
preserves every pairwise order. On trees this condition is exact for a guarantee over all score vectors: any
tree containing a three-edge path admits a reversal. It gives an ordinal guarantee, while P's reward-shaping
use of numerical scores still needs score magnitudes. Ada proved the tree boundary and the diameter-two
extension; Emmy checked the short-path argument.

Cycles can recover more within this restricted model. On a comparison triangle with consistently oriented
calibrated means \(A=w_{ij}\), \(B=w_{jk}\), and \(C=w_{ik}\), the Bradley--Terry addition identity gives
\[
C=\frac{A+B}{1+AB/b^2},\qquad
b^2=\frac{ABC}{A+B-C}.
\]
When the three means are nonzero, the positive admissible root identifies \(b\). Each edge then yields its
score difference, and a connected graph recovers all scores up to a common shift. A nondegenerate four-cycle
also identifies \(b\): if \(S_1\) is the sum of its consistently oriented edge means and \(S_3\) is the sum of
their three-way products, cycle consistency gives \(S_1b^2+S_3=0\), so \(S_1\ne0\) fixes the positive
admissible rate. Emmy derived these identities and Ada checked them. A cycle of arbitrary length is not
enough: Ada supplied a six-cycle with two admissible rates, about \(0.739\) and \(0.854\), whose identical
report laws reverse one vertex pair's order.

\section{Objections and limits}

The four-vertex example objected to treating a fitted agreement matrix as evidence that reports track the
intended comparison. The objection stands: more samples of the same reports cannot distinguish its two worlds.
A direct comparison across that path creates a cycle that exposes this particular clean-model inconsistency,
and a nondegenerate short cycle calibrates \(b\) within the constant-rate model. Neither observation validates
that model against edge-dependent or unrestricted shared bias. Degenerate cycles can also provide no rate
information.

The clean sampling bound answers a different objection, namely that population identification says little
about limited data or query cost. It makes both depend on the observer and comparison graphs, but its
misspecification term need not shrink with more data. No finite-sample stability result for the cycle
estimate, optimal audit design for general graphs, truthful bootstrap for peer reports, or sensitivity theorem
for a particular downstream decision was established. The peers also did not establish that the graph results
are new. Their targeted checks found work on correlated annotators \cite{li2019}, pairwise co-occurrence
identification \cite{ibrahim2019}, semi-verified learning \cite{meister2018}, worker-dependent ranking errors
\cite{crowdbt}, and Bradley--Terry diagnostics \cite{btdiagnostics}. Those checks were limited searches, not a
full review.

\section{Why the thread stopped}

The verifier marked the thread DRAFT because the note's central claims were supported and the pair yielded a
non-obvious result: apparently sound peer agreements can coexist with a wrong downstream ranking. The
graph-dependent order result and conditional cycle calibration made the proposed prediction interface
concrete, while the assumptions and limits remained explicit. The smallest useful restart is to prove a
finite-sample bound for an informative triangle or four-cycle with a stated separation from degeneracy,
counting both comparison labels and graph conditioning. That bound could then be passed to one specified
decision algorithm with an explicit sensitivity rule. The present material does not supply that final step.

\bibliographystyle{plainurl}
\bibliography{references}

\end{document}
